% RESEARCH_DRAFT_SHA256: bda3a83ae1ef8ccbd33972757a05fae4266cf783f6455c05d9a3c2227df9168f
% Layout and source-reference wrappers expanded for the web renderer.
\documentclass[11pt]{article}
\usepackage[margin=25mm]{geometry}
\usepackage{amsthm}
\usepackage{fontspec,unicode-math}
\setmainfont{Latin Modern Roman}
\setmathfont{Latin Modern Math}
\usepackage{amsmath,amssymb,mathtools}
\usepackage{xurl,hyperref}
\hypersetup{colorlinks=true,urlcolor=blue,linkcolor=blue}
\setcounter{secnumdepth}{0}
\setlength{\parindent}{0pt}
\setlength{\parskip}{5pt}
\setlength{\emergencystretch}{3em}
\newcommand{\sref}[2]{\hyperlink{source:#1:#2}{[S#2]}}
\newcommand{\cataloguescope}[1]{\paragraph{Source question.}#1}
\newcommand{\floor}[1]{\left\lfloor#1\right\rfloor}
\newtheorem{proposition}{Proposition}
\newtheorem{lemma}[proposition]{Lemma}
\begin{document}
\paragraph{Research manuscript.} OpenAI \texttt{gpt-6-astra}, reasoning \texttt{ultra}; 11 October 2026.
\paragraph{Preserved source body.} The complete catalogue mathematical text and sources follow with layout and source links expanded, including source qualifications and its historical preparation notes. New research begins at the section headed Research attempt.

\section{Fixed-Parameter Tractability of Zonotope Containment}\label{statement}

\subsection{Definitions and mathematical statement}

\paragraph{Shared notation.}Standard number systems are denoted by $\mathbb N,\mathbb Z,\mathbb Q,\mathbb R,\mathbb C$; index ranges, computational models and specialized notation follow the cited problem source and the definitions below.


The centered zonotope generated by $G$ is $Z(G)=\{Gx:x\in[-1,1]^n\}$. Containment means every point of the first zonotope belongs to the second. Dimension-parameterized runtime has a computable factor depending only on $d$.

Fixed-parameter tractability for parameter $k$ means running time $f(k)N^{O(1)}$, where $N$ is the input length and $f$ is computable. The exponent is independent of $k$.
\paragraph{Statement recorded in the supplied source.}
Given $G\in\mathbb R^{d\times n}$ and $H\in\mathbb R^{d\times m}$, can one decide whether the zonotope $Z(G)$ is contained in $Z(H)$ in time $f(d)\operatorname{poly}(n+m)$ for some computable function $f$? \href{https://doi.org/10.4171/owr/2024/50}{[S2]}
\paragraph{Scope qualification.}
The source\textquotesingle s centered/noncentered zonotope convention and real-input computational model are not specified in the short statement.
\paragraph{Source status.}
The supplied catalogue labels this record open. Its source-status metadata was checked against the supplied snapshot on 2026-10-11. This import does not independently establish current open status or mathematical correctness. The catalogue literature assessment is dated 2026-08-22.
\subsection{Short English statement}
Decide zonotope containment with dimension confined to the parameter-dependent runtime factor.
\subsection{Sources}
\begin{itemize}
\item[S1] Ulam AI and the UnsolvedMath Contributors, UnsolvedMath: A Curated Collection of Open Mathematics Problems, catalogue entry. Catalogue curation and metadata: CC BY 4.0. The supplied source question is retained; definitions and the English summary are editorial additions.
\url{https://huggingface.co/datasets/ulamai/UnsolvedMath}.
\item[S2] Combinatorial Optimization (2025), Oberwolfach Reports, DOI 10.4171/owr/2024/50, pp. 2949--3000
\url{https://doi.org/10.4171/owr/2024/50}.
\item[S3] Kulmburg and Althoff, Zonotope containment
\url{https://www.sciencedirect.com/science/article/pii/S0947358021000856}.
\item[S4] Froese et al., Open problem on zonotope FPT
\url{https://mlanthology.org/colt/2025/froese2025colt-open/}.
\item[S5] Skutella in Oberwolfach Report 50
\url{https://doi.org/10.4171/OWR/2024/50}.
\textit{Underlying Oberwolfach Reports material retains its source-specific licence, including CC BY-SA 4.0 where applicable; see THIRD\_PARTY\_NOTICES.md.}
\end{itemize}
\paragraph{Preparation and review.}
Original question attributed to the cited source and UnsolvedMath contributors. Codex (GPT-6; exact serving revision not exposed) assisted with definitions, English summaries, source inspection and formatting. No new human mathematical authorship or independent proof verification is asserted. Catalogue identity checked against the complete supplied registry and existing repository statements; known flagged repetitions and obvious identical conjectures excluded. Full mathematical-equivalence/literature review has not been performed. Statement preparation and final review followed contribution prompts C and D; automated checks verify structure and identity, not open status.
\label{end:statement}

\section{Research attempt}
\subsection{Literature, source context and scope}
The full catalogue body, including qualifications and historical source status, is preserved above. The original workshop report was inspected in the locally cached publisher PDF text. This attempt distinguishes the source question from any explicitly restricted model introduced below. Repeated records retain their own source IDs and do not become independent conjectures. Elementary results and attributed literature results are not claimed as new original-target progress.

\subsection{Literature qualification and a support-test route}
The source report uses $Z_0(G)=\sum_i[0,g_i]$, while the catalogue defines the centered $Z_c(G)=G[-1,1]^n$. They satisfy
\[
 Z_0(G)=\tfrac12\sum_i g_i+\tfrac12 Z_c(G).
\]
The center displacement matters in containment and cannot be discarded. The later primary paper cited below establishes W[1]-hardness of the source non-containment problem, so a general FPT containment algorithm in its rational-input model would imply $\mathrm{FPT}=\mathrm{W[1]}$. This is attributed literature, not a newly proved resolution here; the original metadata remain preserved.
Our attempted simplification is to certify containment from the $2d$ coordinate support inequalities. For the centered convention take
\[
 G=\binom11,\qquad H=\begin{pmatrix}1/2&1/2\\1/2&-1/2\end{pmatrix}.
\]
\begin{proposition}
The two coordinate projections of $Z_c(G)$ and $Z_c(H)$ agree, but $Z_c(G)\nsubseteq Z_c(H)$.
\end{proposition}
\begin{proof}
Both coordinate ranges are $[-1,1]$. A point of $Z_c(H)$ has coordinates $((s+t)/2,(s-t)/2)$ with $|s|,|t|\le1$, equivalently $|x+y|\le1$ and $|x-y|\le1$. The point $(1,1)\in Z_c(G)$ violates the first inequality. In support notation, maximizing a linear form over independent segment coefficients gives $h_{Z_c(G)}(u)=\sum_i|u^Tg_i|$; hence at $u=(1,1)$ the two supports are two and one.
\end{proof}
This also produces a source-convention example: replace each generator by the pair consisting of it and its negative. Then the noncentered sum of these paired segments equals the original centered zonotope, with center zero. Thus the coordinate-test failure is not an artifact of the convention.
\subsection{Verification and remaining gaps}
The four coordinate supports and the diagonal witness were evaluated exactly. Testing all directions, avoiding dimension-dependent enumeration, and matching the precise input model remain unresolved in this attempt. The published conditional obstruction must not be recast as our unconditional impossibility theorem.

\section{References and source attribution}
The primary Oberwolfach workshop report, \url{https://doi.org/10.4171/owr/2024/50}, was inspected on 11 October 2026. Relevant source passages: 2024/50 pp. 2988--2989, noncentered zonotope definition. Catalogue and contributor attribution remain in the original Sources section. Mathematical arguments labeled as this attempt are proved above; cited external results remain attributed to their authors and are not independently reproved or claimed as our discoveries.

\url{https://arxiv.org/html/2509.22849v3}. Froese, Grillo, Hertrich and Stargalla, Parameterized Hardness of Zonotope Containment and Neural Network Verification, Theorem 6.1: non-containment is W[1]-hard in dimension. The theorem and support-function reduction were inspected. This is an external conditional complexity obstruction, not our proof; a containment FPT algorithm would also decide its complement in FPT. The preserved dated catalogue assessment does not incorporate this attributed hardness result.


\paragraph{Integration review.} Codex (GPT-6; exact serving revision not exposed) reviewed the written arguments and expanded layout and source-reference wrappers. Repository metadata and the production web renderer were checked using Python 3.13 and JavaScript with MathJax 3.2.2. This model-assisted review is not independent human proof verification; no PDF compilation or formal proof checking is asserted.

\section{Completion Estimate}
\noindent\textbf{COMPLETION: 0\%}
No substantiated new progress toward the entire original target is claimed. The worked results and failed route are delimited above; a known result, illustrative model, or auxiliary obstruction is not counted as original-target completion.

\section{Final status and attribution}
\textbf{Attempt status: unresolved.} This describes this research attempt, not a new ruling on the source catalogue's status or on external literature claims.
The mathematics and this self-audit were generated by OpenAI \texttt{gpt-6-astra}, reasoning \texttt{ultra}, on 11 October 2026. Codex (GPT-6; exact serving revision not exposed) orchestrated the batch and contributed formatting/repository review. No human mathematical authorship, independent human proof checking, or formal verification is claimed. The historical preparation paragraphs in the reproduced catalogue body concern that body only.

Contribution prompt A was applied to the complete definition and supplied source record, followed by the adversarial checks of prompt B. Local reproduction files are \texttt{batch011/checks.py}, \texttt{computations.json}, and the per-ID audit JSON. They were not public links at preparation time. Numerical checks are finite and do not prove unrestricted results. No LaTeX engine was found locally; compilation was not performed by this research agent.

\end{document}
